\rhead{MA202: Linear Algebra \& Opt.}
\phantomsection 
\section*{Linear Algebra \& Optimization (MA202)}
\label{course:MA202}

\noindent \small{\hyperref[main:scheme]{$\leftarrow$ Back to Scheme}} \vspace{0.5cm}

\noindent \textbf{Credits:} 3 (3-0-0)

\subsection*{Course Content}
\noindent\textbf{Unit 1} \\
\textit{Matrices and Linear Systems:} Matrix Algebra, Rank of a matrix, Reduction to Echelon and Normal forms. Consistency and solution of Systems of Linear Equations (Homogeneous and Non-homogeneous) using Gauss Elimination and Gauss-Jordan methods. LU Decomposition (Algorithmic approach to solving systems). Inverse of a matrix via elementary row operations.

\vspace{0.3cm}

\noindent\textbf{Unit 2} \\
\textit{Vector Spaces and Inner Products:} Vector Spaces: Subspaces, Linear Independence/Dependence, Span, Basis, and Dimension. Linear Transformations: Kernel, Range, Rank-Nullity Theorem, and Matrix representation of transformations. Inner Product Spaces: Norms, Orthogonality, Orthogonal Projections, and the Gram-Schmidt Orthogonalization Process (QR Decomposition foundation).

\vspace{0.3cm}

\noindent\textbf{Unit 3} \\
\textit{Eigen Theory and Matrix Factorization:} Eigenvalues and Eigenvectors: Characteristic equation, Algebraic and Geometric Multiplicity. Cayley-Hamilton Theorem. Diagonalization of Matrices. Real Symmetric Matrices and Orthogonal Diagonalization. Advanced Factorization: Introduction to Singular Value Decomposition (SVD) and Principal Component Analysis (PCA) (Mathematical foundation for Data Compression and Dimensionality Reduction).

\vspace{0.3cm}

\noindent\textbf{Unit 4} \\
\textit{Linear Programming (Optimization):} Introduction to Optimization and Convex Sets. Formulation of Linear Programming Problems (LPP). Solution methods: Graphical Method and the Simplex Algorithm (Big-M method, Two-phase method). Duality in Linear Programming: Primal-Dual relationships and Sensitivity Analysis (Economic interpretation of resources).

\vspace{0.3cm}

\noindent\textbf{Unit 5} \\
\textit{Non-Linear and Numerical Optimization:} Unconstrained Optimization: Necessary and sufficient conditions for local extrema. Iterative Search Methods: Gradient Descent (Steepest Descent), Newton-Raphson Method, and Conjugate Gradient Method (Training logic for Neural Networks). Constrained Optimization: Lagrange Multipliers (Revisited) and Karush-Kuhn-Tucker (KKT) conditions for non-linear constraints.

\newpage
\rhead{Curriculum Schema}